%!TEX root = ../sztuthesis_main.tex
% 论文正文是主体，主体部分应从另页右页开始，每一章应另起页。一般由序号标题、文字叙述、图、表格和公式等五个部分构成。
\section{引言}
\subsection{研究背景与意义}

\app{Word} 不难用，但是想用得漂亮还得费一番功夫。

对于小白来说，（注意！！是对于小白来说，不要跟我杠！！！我就是 \app{Word} 小白，高端玩法玩不动）：

插入个图片，下面的说明文字是不是插入文本框？那文本框要不要跟图片“组合”？是不是直接圈没法圈起来？因为图片要变成浮动格式，和文本框绑定后再改回嵌入格式，你说蛋疼不蛋疼？不组合？那有一定概率发生你的图片在上一页，描述文字在下一页。呵呵。

插入参考文献，手动编辑？我的天哪，一百多个文献，中间插一个，怎么改序号？
很好，可以用交叉引用，一个个编辑文献格式？
很好，可以用 \app{endnote} 或者 \app{noteexpress} 的插入功能，你有没有发现插入是个宏，\oper{ctrl+z} 的时候会烦死你啊……
叮叮！让我们祭出 \LaTeX！！\cite{2013浅谈使用}，有 \env{bibliography}，一个 \pkg{cite} 包打天下！不要太爽。


\begin{figure}[hbt]
    \centering
    \includegraphics[width=0.5\textwidth]{latex.jpg}
    \caption{LaTeX}
    \label{F.latex}
\end{figure}

插入公式，对 \app{Word} 小白来说，公式居中编号靠右就是一道百度搜索能力过滤器。

\app{Word} 里编辑三线表，啊烦躁。

等等等等……

让我们，专心写论文好不好？

爱你们。


\subsection{主要研究工作}
虽然我 \LaTeX 水平也很水……但是通过大量 \oper{debug} 也勉强给大家凑出来一个格式绝对标准的 \LaTeX 模板，模板代码丑就丑吧，能用就行。写了大量注释，有一点 \LaTeX 基础就可以根据自己需要修改 \cls{SZTUthesis.cls} 文件。

(1) 提供图片插入示例。

(2) 提供表格插入示例。

(3) 提供公式插入示例。

(4) 提供参考文献插入示例。

\subsection{论文组织结构}

全文内容共六章，具体内容组织如下：

第一章为绪论。

第二章为图片插入示例。

第三章为表格插入示例。

第四章为公式插入示例。

第五章为参考文献插入示例。

第六章总结与展望，总结了本文的主要工作，展望了下一阶段的研究方向。

% \newpage    % 两个章节之间分页，不想分的话可注释掉

\section{图像布局}
\label{sec.figure}

\emph{学校对图片只有小标题要求，没有进一步的子图要求，我们按科技论文常规排版来}

\subsection{单图布局}

\lipsum

\emph{单图布局如图~\ref{F.sztu_single} 所示。}

\begin{figure}[hbt]
\centering
\includegraphics[width=0.5\textwidth]{sztu.png}
\caption{单图布局示例}
\label{F.sztu_single}
\end{figure}

\subsection{横排布局}

\emph{横排布局如图~\ref{F.sztu_row} 所示。}

\begin{figure}[!htb]
    \centering
    \begin{subfigure}[t]{0.24\linewidth}
        \begin{minipage}[b]{1\linewidth}
        \includegraphics[width=1\linewidth]{sztu.png}
        \caption{可以增加描述}
        \end{minipage}
    \end{subfigure}
    \begin{subfigure}[t]{0.24\linewidth}
        \begin{minipage}[b]{1\linewidth}
        \includegraphics[width=1\linewidth]{sztu.png}
        \caption{}
        \end{minipage}
    \end{subfigure}
    \begin{subfigure}[t]{0.24\linewidth}
        \begin{minipage}[b]{1\linewidth}
        \includegraphics[width=1\linewidth]{sztu.png}
        \caption{}
        \end{minipage}
    \end{subfigure}
    \begin{subfigure}[t]{0.24\linewidth}
        \begin{minipage}[b]{1\linewidth}
        \includegraphics[width=1\linewidth]{sztu.png}
        \caption{}
        \end{minipage}
    \end{subfigure}
    \caption{横排布局示例}
    \label{F.sztu_row}
\end{figure}

\lipsum

\subsection{竖排布局}
\emph{竖排布局如图\ref{F.sztu_col}所示。}

\begin{figure}[!htb]
    \centering
    \begin{subfigure}[t]{0.15\linewidth}
        \begin{minipage}[b]{1\linewidth}
        \includegraphics[width=1\linewidth]{sztu.png}
        \caption{}
        \end{minipage}
    \end{subfigure}\\
    \begin{subfigure}[t]{0.15\linewidth}
        \begin{minipage}[b]{1\linewidth}
        \includegraphics[width=1\linewidth]{sztu.png}
        \caption{}
        \end{minipage}
    \end{subfigure}
    \caption{竖排布局示例}
    \label{F.sztu_col}
\end{figure}

\lipsum

\subsection{竖排多图横排布局}

\begin{figure}[!htb]
    \centering
    \begin{subfigure}[t]{0.13\linewidth}
        \begin{minipage}[b]{1\linewidth}
        \includegraphics[width=1\linewidth]{sztu.png} \vspace{-1ex} \vfill
        \includegraphics[width=1\linewidth]{sztu.png}
        \end{minipage}
        \caption{}
    \end{subfigure}
    \begin{subfigure}[t]{0.13\linewidth}
        \begin{minipage}[b]{1\linewidth}
        \includegraphics[width=1\linewidth]{sztu.png} \vspace{-1ex} \vfill
        \includegraphics[width=1\linewidth]{sztu.png}
        \end{minipage}
        \caption{}
    \end{subfigure}
    \caption{竖排多图横排布局}
    \label{F.sztu_col_row}
\end{figure}

\emph{竖排多图横排布局如图~\ref{F.sztu_col_row} 所示。注意看(a)、(b)编号与图关系。}


\subsection{横排多图竖排布局}

\lipsum

\begin{figure}[!htb]
    \centering
    \begin{subfigure}[t]{0.3\linewidth}
        \begin{minipage}[b]{1\linewidth}
        \includegraphics[width=0.45\linewidth]{sztu.png}
        \includegraphics[width=0.45\linewidth]{sztu.png}
        \end{minipage}
        \caption{}
    \end{subfigure}\\
    \begin{subfigure}[t]{0.3\linewidth}
        \begin{minipage}[b]{1\linewidth}
        \includegraphics[width=0.45\linewidth]{sztu.png}
        \includegraphics[width=0.45\linewidth]{sztu.png}
        \end{minipage}
        \caption{}
    \end{subfigure}
    \caption{横排多图竖排布局}
    \label{F.sztu_row_col}
\end{figure}

\emph{横排多图竖排布局如图~\ref{F.sztu_row_col} 所示。注意看(a)、(b)编号与图关系。}

\subsection{本章小结}
本章示例图片布局。

% \newpage    % 两个章节之间分页，不想分的话可注释掉


\section{表格插入示例}

\begin{table}[htb]
  \centering
  \caption{学校文件里对表格的要求不是很高，不过按照学术论文的一般规范，表格为三线表。}
  \label{T.example}
  \begin{tabular}{llllll}
  \hline
   & A  & B  & C  & D  & E \\
  \hline
1 	& 212 & 414 & 4 		& 23 & fgw	\\
2 	& 212 & 414 & v 		& 23 & fgw	\\
3 	& 212 & 414 & vfwe		& 23 & 嗯	\\
4 	& 212 & 414 & 4fwe		& 23 & 嗯	\\
5 	& af2 & 4vx & 4 		& 23 & fgw	\\
6 	& af2 & 4vx & 4 		& 23 & fgw	\\
7 	& 212 & 414 & 4 		& 23 & fgw	\\

\hline{}
\end{tabular}
\end{table}

\emph{表格如表~\ref{T.example} 所示，\LaTeX 表格技巧很多，这里不再详细介绍。}

\lipsum

% \newpage    % 两个章节之间分页，不想分的话可注释掉

\section{公式插入示例}

\lipsum

\emph{公式插入示例如公式~\eqref{E.example} 所示。}
\begin{equation}
\gamma_x=
\begin{cases}
  0, & \text{if $|x| \leq \delta$} \\
  x, & \text{otherwise}
\end{cases}
\label{E.example}
\end{equation}


% \newpage    % 两个章节之间分页，不想分的话可注释掉

\section{参考文献插入示例}

\LaTeX \cite{lamport1994latex}插入参考文献最方便的方式是使用 \env{bibliography}\cite{pritchard1969statistical}，大多数出版商的论文页面都会有导出 \format{bib} 格式参考文献的链接，把每个文献的 \format{bib} 放入 \bib{thesis-references.bib}，然后用 \oper{bibkey} 即可插入参考文献。

\lipsum

% \newpage    % 两个章节之间分页，不想分的话可注释掉


\section{总结与展望}

\noindent{纯数字编号}
\begin{enumerate}
 \item XXXXXXXXXX
 \label{item1}
 \item XXXXXXXXXX
 \item XXXXXXXXXX
\end{enumerate}
罗马编号
\begin{enumerate}[label=(\roman*)]
 \item XXXXXXXXXX
 \label{item2}
 \item XXXXXXXXXX
 \item XXXXXXXXXX
\end{enumerate}
括号编号
\begin{enumerate}[label=(\arabic*)]
 \item XXXXXXXXXX
 \label{item3}
 \item XXXXXXXXXX
 \item XXXXXXXXXX
\end{enumerate}
半括号编号
\begin{enumerate}[label=\arabic*)]
 \item XXXXXXXXXX
 \label{item4}
 \item XXXXXXXXXX
 \item XXXXXXXXXX
\end{enumerate}
小字母编号
\begin{enumerate}[label=\alph*)]
 \item XXXXXXXXXX
 \label{item5}
 \item XXXXXXXXXX
 \item XXXXXXXXXX
\end{enumerate}

引用测试,正如~\ref{item1}、\ref{item2}、\ref{item3}、\ref{item4}、\ref{item5} 所示

\subsection{工作展望}
手动编号 %（不推荐,无法被交叉引用）

本课题针对XX，鉴于XXX，对XX进行了提高，但是XXX，所以有如下XX：

（1）目前XX虽然XX，但是XX仍然XX，所以XX仍然是一个值得XX的问题。

（2）随着XX，XX具有XX的问题，仍值得进一步XX。

（3）本课题在XX有了XX，但是XX的XX还存在XX，所以XX。


\newpage
